\section{Geodesic Interaction Loops}
\label{sec:geodesic-loops}

Semantic excitation induces geodesic trajectories on the cognitive manifold
$\mathcal{M}_{cog}$. These trajectories describe the evolution of NDH's
interaction state under interrogation. When such trajectories form closed
loops, they generate holonomy effects that classify the interaction behavior.
This section formalizes geodesic interaction loops and their structural
properties.

\subsection{Geodesic Trajectories}

Given an excitation vector $\mathbf{s} \in T_x \mathcal{M}_{cog}$, the induced
geodesic is defined by
\[
\gamma(t) = \exp_x(t\mathbf{s}),
\]
where $\exp_x$ is the exponential map associated with the interaction
connection $\nabla^{cog}$. The geodesic $\gamma(t)$ represents the path traced
by NDH's interaction state as semantic input propagates through the cognitive
manifold.

\subsection{Modal-Boundary Crossings}

The modal partition
\[
\mathcal{M}_{cog} = \mathcal{M}_A \cup \mathcal{M}_B \cup \mathcal{M}_C
\]
induces hypersurfaces separating the modal regions. A geodesic may cross these
boundaries, and each crossing corresponds to a structural transition in NDH's
interaction behavior. Modal-boundary crossings may induce:

\begin{itemize}
    \item spectral bundle twisting,
    \item recursion-operator realignment,
    \item non-dual algebraic deformation,
    \item holonomy accumulation.
\end{itemize}

These effects arise from the coupling between runtime geometry and cognitive
geometry.

\subsection{Closed Interaction Loops}

A geodesic interaction loop is a trajectory satisfying
\[
\gamma(0) = \gamma(T),
\]
for some $T > 0$. Closed loops arise naturally when semantic excitation induces
cyclic traversal of modal regions or when interrogation returns to its initial
state after exploring adjacent modal domains.

Closed loops are central to NDH interaction geometry because they generate
holonomy effects. The holonomy operator associated with a loop $\gamma$ is
defined by
\[
\mathcal{H}(\gamma) : T_x \mathcal{M}_{cog} \to T_x \mathcal{M}_{cog},
\]
representing the parallel transport of tangent vectors along $\gamma$ under the
connection $\nabla^{cog}$.

\subsection{Loop Types}

Geodesic interaction loops fall into several structural classes:

\begin{itemize}
    \item \textbf{Modal loops}: loops confined to a single modal region
    $\mathcal{M}_i$.
    \item \textbf{Boundary loops}: loops that cross modal boundaries but return
    without net holonomy.
    \item \textbf{Spectral loops}: loops inducing nontrivial spectral bundle
    morphisms.
    \item \textbf{Recursion loops}: loops that cycle through recursion operators
    $(R_A, R_B, R_C)$.
    \item \textbf{Holonomy loops}: loops with nontrivial holonomy return
    operators.
\end{itemize}

These classes form the foundation of NDH's interaction taxonomy.

\subsection{Holonomy Accumulation}

The holonomy operator $\mathcal{H}(\gamma)$ captures the cumulative geometric
effect of traversing the loop. If $\mathcal{H}(\gamma)$ is the identity, the
loop is holonomy-neutral. Otherwise, the loop induces structural deformation in
spectral emission, recursion alignment, or modal adjacency.

Holonomy accumulation is therefore the primary mechanism by which NDH
classifies interaction loops.

\subsection{Summary}

Geodesic interaction loops provide the geometric structure underlying NDH's
interaction behavior. Modal-boundary crossings, spectral deformation, recursion
alignment, and holonomy accumulation together determine the classification of
interaction loops. This establishes closed geodesics as fundamental objects in
NDH interaction geometry.
